\chapter{Ultrailoczyny i twierdzenie Łosia}
\label{ch:mod-ultraproducts}

\section{Motywacja i ultrafiltry}
Iloczyn prosty struktur nie zachowuje na ogół wszystkich zdań
pierwszego rzędu prawdziwych w każdym czynniku. Chcemy skonstruować
iloraz iloczynu, w którym prawdziwość formuły zależy od tego, czy
zbiór indeksów, dla których jest ona prawdziwa, jest wystarczająco duży.
Ultrafiltr rozstrzyga tę kwestię dla każdego podzbioru indeksów.
W całym rozdziale pracujemy w ZFC i rozważamy struktury o niepustych
uniwersach dla jednej ustalonej sygnatury pierwszego rzędu z równością,
w sensie definicji \ref{def:sem-structure}.

Przypomnijmy jedynie potrzebne własności: ultrafiltr $\mathcal U$
na niepustym zbiorze $I$ zawiera $I$, nie zawiera $\varnothing$,
jest zamknięty na skończone przecięcia i przechodzenie do nadzbiorów,
a dla każdego $A\subseteq I$ dokładnie jeden ze zbiorów $A,I\setminus A$
należy do $\mathcal U$. Definicję i kryterium maksymalności
omówiono w teorii mnogości (definicja \ref{def:bpi-boolean-filter}
i lemat \ref{lem:bpi-ultrafilter-criterion}).
W szczególności $A\cap B\in\mathcal U$ wtedy i tylko wtedy, gdy
$A,B\in\mathcal U$. Jeśli $A\cap D=B\cap D$ dla pewnego $D\in\mathcal U$,
to $A\in\mathcal U$ wtedy i tylko wtedy, gdy $B\in\mathcal U$:
wynika to z zamknięcia na przecięcia i nadzbiory.

\section{Konstrukcja ultrailoczynu}
\label{sec:mod-ultraproduct-construction}
Niech $(\mathfrak M_i)_{i\in I}$ będzie rodziną struktur o uniwersach
$M_i$. Aksjomat wyboru zapewnia, że $P=\prod_{i\in I}M_i$
jest niepusty.

\begin{definicja}[Równość modulo ultrafiltr]
\label{def:mod-ultra-equivalence}
Dla $a,b\in P$ kładziemy
\[
a\sim_{\mathcal U}b
\quad\Longleftrightarrow\quad
\{i\in I:a_i=b_i\}\in\mathcal U.
\]
Klasę równoważności $a$ oznaczamy przez $[a]_{\mathcal U}$.
\end{definicja}

\begin{lemat}
\label{lem:mod-ultra-equivalence}
Relacja $\sim_{\mathcal U}$ jest relacją równoważności.
\end{lemat}
\begin{proof}
Zwrotność wynika z $I\in\mathcal U$, a symetria z symetrii równości.
Jeżeli $a\sim_{\mathcal U}b$ i $b\sim_{\mathcal U}c$, to przecięcie
zbiorów indeksów, na których zachodzą obie równości, należy do
$\mathcal U$ i jest podzbiorem $\{i:a_i=c_i\}$. Stąd przechodniość.
\end{proof}

\begin{definicja}[Ultrailoczyn i ultrapotęga]
\label{def:mod-ultraproduct}
Ultrailoczyn $\mathfrak M=\prod_{i\in I}\mathfrak M_i/\mathcal U$
ma uniwersum $P/{\sim_{\mathcal U}}$. Dla stałej $c$, symbolu
funkcyjnego $f$ arności $n\geq 1$ oraz symbolu relacyjnego $R$
arności $n$ określamy
\begin{align*}
c^{\mathfrak M}&=[(c^{\mathfrak M_i})_{i\in I}]_{\mathcal U},\\
f^{\mathfrak M}([a^1]_{\mathcal U},\ldots,[a^n]_{\mathcal U})
&=[(f^{\mathfrak M_i}(a_i^1,\ldots,a_i^n))_{i\in I}]_{\mathcal U},\\
\mathfrak M\models R([a^1]_{\mathcal U},\ldots,[a^n]_{\mathcal U})
&\Longleftrightarrow
\{i:\mathfrak M_i\models R(a_i^1,\ldots,a_i^n)\}\in\mathcal U.
\end{align*}
Dla symbolu relacyjnego arności zero ostatnia reguła dotyczy
zdania atomowego bez argumentów. Gdy wszystkie czynniki są
równe $\mathfrak A$, otrzymujemy ultrapotęgę $\mathfrak A^I/\mathcal U$.
\end{definicja}

\begin{lemat}[Niezależność od reprezentantów]
\label{lem:mod-ultra-welldefined}
Interpretacje symboli w definicji \ref{def:mod-ultraproduct}
są dobrze określone.
\end{lemat}
\begin{proof}
Jeżeli $[a^k]_{\mathcal U}=[b^k]_{\mathcal U}$ dla $1\leq k\leq n$,
to $D=\bigcap_{k=1}^n\{i:a_i^k=b_i^k\}\in\mathcal U$.
Wartości funkcji zgadzają się na $D$, więc wyznaczają tę samą
klasę. Zbiory prawdziwości relacji mają równe przecięcia z $D$,
a zatem jednocześnie należą lub nie należą do ultrafiltru.
Stałe i relacje arności zero nie wymagają wyboru reprezentantów.
\end{proof}

\begin{lemat}[Wartość termu]
\label{lem:mod-ultra-terms}
Dla termu $t(x_1,\ldots,x_n)$ zachodzi
\[
t^{\mathfrak M}([a^1]_{\mathcal U},\ldots,[a^n]_{\mathcal U})
=[(t^{\mathfrak M_i}(a_i^1,\ldots,a_i^n))_{i\in I}]_{\mathcal U}.
\]
\end{lemat}
\begin{proof}
Indukcja po budowie termu. Dla zmiennej obie strony są jej klasą,
dla stałej wynika to z definicji. Jeśli $t=f(t_1,\ldots,t_m)$,
stosujemy hipotezę indukcyjną do każdego $t_j$, a następnie
definicję interpretacji $f$ w ultrailoczynie.
\end{proof}

\section{Twierdzenie Łosia}
\label{sec:mod-los}
\begin{twierdzenie}[Twierdzenie Łosia]
\label{thm:mod-los}
Dla każdej formuły pierwszego rzędu $\varphi(\bar x)$,
której zmienne wolne należą do skończonej listy $\bar x$,
i dowolnych reprezentantów $\bar a=(a^1,\ldots,a^n)$ mamy
\[
\mathfrak M\models\varphi([\bar a]_{\mathcal U})
\quad\Longleftrightarrow\quad
\{i\in I:\mathfrak M_i\models\varphi(\bar a_i)\}\in\mathcal U.
\]
Zapis $[\bar a]_{\mathcal U}$ oznacza krotkę klas.
\end{twierdzenie}
\begin{proof}
Piszmy $S_\psi(\bar a)=\{i:\mathfrak M_i\models\psi(\bar a_i)\}$.
Dowód przeprowadzimy przez indukcję po budowie formuły.
W przypadku atomowej równości $t=s$ korzystamy z lematu
\ref{lem:mod-ultra-terms} i definicji równości klas.
Dla atomowej relacji $R(t_1,\ldots,t_m)$ ten sam lemat i
definicja interpretacji $R$ dają żądaną równoważność.

Dla negacji $S_{\neg\psi}=I\setminus S_\psi$, więc
$S_{\neg\psi}\in\mathcal U$ dokładnie wtedy, gdy
$S_\psi\notin\mathcal U$. Dla koniunkcji
$S_{\psi\land\theta}=S_\psi\cap S_\theta$, a przecięcie należy
do ultrafiltru dokładnie wtedy, gdy oba zbiory do niego należą.
Hipoteza indukcyjna daje w obu przypadkach tezę.

Pozostaje kwantyfikator egzystencjalny. Niech
$S=\{i:\mathfrak M_i\models\exists y\,\psi(y,\bar a_i)\}$.
Jeśli $\mathfrak M\models\exists y\,\psi(y,[\bar a]_{\mathcal U})$,
wybieramy świadka $[b]_{\mathcal U}$ w ultrailoczynie.
Hipoteza indukcyjna daje
$\{i:\mathfrak M_i\models\psi(b_i,\bar a_i)\}\in\mathcal U$.
Zbiór ten zawiera się w $S$, zatem $S\in\mathcal U$.

Odwrotnie, jeśli $S\in\mathcal U$, dla $i\in S$ wybieramy
$b_i\in M_i$ spełniające $\psi(b_i,\bar a_i)$, a dla $i\notin S$
dowolny element $M_i$. Korzystamy tu z aksjomatu wyboru.
Zbiór $\{i:\mathfrak M_i\models\psi(b_i,\bar a_i)\}$
zawiera $S$, więc należy do $\mathcal U$. Z hipotezy indukcyjnej
wynika $\mathfrak M\models\psi([b]_{\mathcal U},[\bar a]_{\mathcal U})$,
co daje formułę egzystencjalną.

Pozostałe spójniki i kwantyfikator ogólny wyrażamy za pomocą
negacji, koniunkcji i kwantyfikatora egzystencjalnego.
\end{proof}

\begin{wniosek}[Zachowanie teorii]
\label{cor:mod-ultra-theory}
Jeśli $\mathfrak M_i\models T$ dla wszystkich $i$ należących
do pewnego zbioru $J\in\mathcal U$, to ultrailoczyn jest modelem $T$.
Wystarczy nawet, aby dla każdego $\sigma\in T$ zbiór
$\{i:\mathfrak M_i\models\sigma\}$ należał do $\mathcal U$.
\end{wniosek}
\begin{proof}
Stosujemy twierdzenie \ref{thm:mod-los} osobno do każdego zdania
$\sigma\in T$. Nie potrzeba przecięcia po wszystkich zdaniach teorii.
\end{proof}

\begin{wniosek}[Zanurzenie diagonalne]
\label{cor:mod-ultrapower-diagonal}
Odwzorowanie $d:\mathfrak A\to\mathfrak A^I/\mathcal U$,
$d(a)=[(a)_{i\in I}]_{\mathcal U}$, jest zanurzeniem elementarnym.
\end{wniosek}
\begin{proof}
Dla krotki stałych reprezentantów zbiór indeksów prawdziwości
dowolnej formuły jest równy $I$ albo $\varnothing$ zależnie
od jej prawdziwości w $\mathfrak A$. Twierdzenie \ref{thm:mod-los}
daje zachowanie i odbijanie prawdziwości wszystkich formuł,
w szczególności równości, co zapewnia też różnowartościowość.
\end{proof}

\section{Przykłady i granice konstrukcji}
\begin{przyklad}[Ultrafiltr główny]
\label{ex:mod-principal-ultraproduct}
Jeśli $\mathcal U=\{A\subseteq I:i_0\in A\}$, to
$[a]_{\mathcal U}\mapsto a_{i_0}$ jest izomorfizmem
ultrailoczynu z $\mathfrak M_{i_0}$. Równość klas jest
dokładnie równością współrzędnych $i_0$. Odwzorowanie jest
na: pozostałe współrzędne można wypełnić elementami ustalonej
krotki z $P$. Interpretacje symboli są zachowane na mocy
definicji \ref{def:mod-ultraproduct}.
\end{przyklad}

\begin{przyklad}[Ultrapotęga arytmetyki]
\label{ex:mod-ultrapower-naturals}
Niech $\mathcal U$ będzie ultrafiltrem niegłównym na $\N$,
którego istnienie omówiono w przykładzie \ref{ex:bpi-free-ultrafilter}.
W ultrapotędze standardowej struktury arytmetycznej
$\mathfrak N^\N/\mathcal U$ element $h=[(n)_{n\in\N}]_{\mathcal U}$
jest większy od każdego $d(k)$, $k\in\N$. Istotnie zbiór
$\{n:n>k\}$ jest ko-skończony, więc należy do $\mathcal U$:
ultrafiltr niegłówny nie zawiera zbioru skończonego, gdyż
z przynależności skończonej sumy wynikałaby przynależność
jednego singletonu. Twierdzenie \ref{thm:mod-los} daje nierówność.
Zanurzenie $d$ jest elementarne, lecz nie jest na.
\end{przyklad}

\begin{przyklad}[Ciała skończone i charakterystyka zero]
\label{ex:mod-ultraproduct-fields}
Niech $(p_n)_{n\in\N}$ będzie ciągiem parami różnych liczb
pierwszych, a $\mathcal U$ ultrafiltrem niegłównym na $\N$.
Rozważmy $K=\prod_n\mathbb F_{p_n}/\mathcal U$ w języku ciał.
Aksjomaty ciała zachodzą w każdym czynniku, więc zachodzą w $K$.
Dla każdej dodatniej liczby całkowitej $m$ zdanie
$\underbrace{1+\cdots+1}_{m}\ne0$ zachodzi we wszystkich
czynnikach poza skończoną liczbą, ponieważ tylko skończenie
wiele liczb pierwszych dzieli $m$. Na mocy twierdzenia Łosia
zdanie to zachodzi w $K$. Zatem $K$ jest ciałem charakterystyki
zero, chociaż każdy czynnik ma charakterystykę dodatnią.
\end{przyklad}

\begin{przyklad}[Iloczyn prosty nie wystarcza]
\label{ex:mod-direct-product-not-elementary}
W języku pierścieni każde ciało spełnia zdanie
$\forall x\,(x^2=x\Rightarrow (x=0\lor x=1))$, ponieważ
$x(x-1)=0$ i ciało nie ma dzielników zera.
Iloczyn prosty dwóch nietrywialnych ciał nie spełnia tego zdania:
element $(0,1)$ jest idempotentem różnym od $0$ i $1$.
Ultrailoczyn ciał natomiast spełnia je na mocy twierdzenia Łosia.
\end{przyklad}

\section{Drugi dowód twierdzenia o zwartości}
\begin{twierdzenie}
\label{thm:mod-compactness-ultraproduct}
Każda skończenie spełnialna teoria pierwszego rzędu $T$ ma model.
\end{twierdzenie}
\begin{proof}
Niech $I$ będzie zbiorem skończonych podzbiorów $T$.
Dla $\Delta\in I$ wybieramy model $\mathfrak M_\Delta\models\Delta$.
Wybór można wykonać na zbiorze: na mocy dolnego twierdzenia
Löwenheima--Skolema każdy spełnialny fragment ma model
o mocy co najwyżej $\max(\aleph_0,|L|)$, gdzie $L$ jest sygnaturą.
Po przeniesieniu uniwersów na podzbiory ustalonej liczby
kardynalnej tej mocy struktury tworzą zbiór; stosujemy AC.

Dla $\sigma\in T$ niech $I_\sigma=\{\Delta\in I:\sigma\in\Delta\}$.
Rodzina tych zbiorów ma własność skończonych przecięć:
dla $\sigma_1,\ldots,\sigma_m$ ich zbiór należy do przecięcia
$I_{\sigma_1}\cap\cdots\cap I_{\sigma_m}$.
Generuje więc filtr właściwy, który przedłużamy do ultrafiltru
$\mathcal U$ (zasada ultrafiltru, definicja \ref{def:bpi-principles}).
Dla każdego $\sigma\in T$ zbiór indeksów prawdziwości $\sigma$
zawiera $I_\sigma$ i należy do $\mathcal U$. Twierdzenie
\ref{thm:mod-los} pokazuje, że
$\prod_{\Delta\in I}\mathfrak M_\Delta/\mathcal U\models T$.
\end{proof}

\begin{uwaga}
\label{rem:mod-ultra-compactness-scope}
Jest to alternatywne wyprowadzenie twierdzenia \ref{thm:compactness},
a nie dowód niezależny od wszystkich wcześniejszych wyników:
w powyższym doborze rodziny użyliśmy dolnego twierdzenia
Löwenheima--Skolema. Twierdzenie Łosia samo nie zależy od
zwartości. W teorii mnogości, w rozdziale \ref{ch:large-cardinals},
rozważano też ultraprodukt uniwersum $V$; dodatkowe trudności
z klasami właściwymi i dobrym ufundowaniem nie występują w
definicji zbiorowego ultrailoczynu struktur, choć jego relacje
nie muszą być dobrze ufundowane.
\end{uwaga}
